\sectionkicker{Source-backed technical notes}
\subsection*{小型化は縮小版ではない――端末・プライバシー・ローカル実行: Technical Notes}
本欄は記事本文で圧縮した一次資料上の情報を、比較・再検証しやすい形で整理したものである。時系列、確認済みの主張、留意点、一次資料URLを掲載する。
\smallskip
\noindent{\small\textit{Technical Notesでは、一次資料から確認した公開・提供・研究上の事実を記載する。数値・能力評価や提供元・著者による比較表現は、独立再現済みの結果ではなく帰属claimとして扱う。}}\par
\medskip
\subsection*{Theme at a glance}
\begin{center}
\footnotesize
\begin{tabularx}{\linewidth}{@{}p{0.20\linewidth}p{0.10\linewidth}p{0.14\linewidth}X@{}}
\toprule
資料 & 位置づけ & 種別 & 時系列 \\
\midrule
Apple Foundation Models & 主要資料 & モデル & 2024-06-10           \\
Phi-3 family & 主要資料 & オープンウェイト & 2024-04-23, 2024-05-21           \\
MobileLLM & 主要資料 & 論文 & 年表対象日付なし \\
OpenELM & 補足資料 & オープンウェイト & 2024-04-22           \\
Gemma family (Gemma / PaliGemma / Gemma 2) & 補足資料 & オープンウェイト & 2024-02-21, 2024-05-14, 2024-06-27           \\
BitNet b1.58 & 補足資料 & 論文 & 年表対象日付なし \\
\bottomrule
\end{tabularx}
\end{center}
\normalsize

\Needspace{0.38\textheight}
\begin{technicalnote}{Apple Foundation Models}{主要資料}
\begingroup
% reader-facing Technical Notes break policy
\widowpenalty=10000
\clubpenalty=10000
\displaywidowpenalty=10000
\begin{tabularx}{\linewidth}{@{}>{\bfseries}p{0.22\linewidth}X@{}}
組織 & Apple \\
種別 & モデル \\
時系列 & 2024-06-10 \\
\end{tabularx}
\smallskip
\begin{minipage}{\linewidth}
% reader-facing Technical Notes event-only fact/source tail group
{\bfseries 一次資料から整理したtechnical points}
\begin{itemize}[leftmargin=1.5em,itemsep=0.35em]
\item \textbf{一次情報で確認できる事実}: Appleの一次資料により、Apple Foundation Modelsについて2024-06-10にResearch公開を確認した。 AppleがWWDC24で説明したfoundation-model構成は、約3B parametersのon-device language modelと、Apple silicon servers上のPrivate Cloud Computeで動くlarger server modelを組み合わせる。両モデルはgrouped-query attentionを使い、on-device側ではmixed 2-bit/4-bit LoRA-based palletizationを含む低bit最適化を適用するほか、featureごとのadapterをbase modelへ動的にload/swapしてtask specializationを行う。性能値はAppleによる評価として扱う。
\end{itemize}
\begin{minipage}{\linewidth}
% reader-facing Technical Notes generic boundary/source tail group
\begin{samepage}
{\bfseries 一次資料}
\begingroup\sloppy
\Urlmuskip=0mu plus 2mu\relax
\begin{itemize}[leftmargin=1.5em,itemsep=0.25em]
\item {\scriptsize\url{https://machinelearning.apple.com/research/introducing-apple-foundation-models}}
\end{itemize}
\endgroup
\end{samepage}
\end{minipage}
\end{minipage}
\endgroup
\end{technicalnote}
\medskip

\Needspace{0.38\textheight}
\begin{technicalnote}{Phi-3 family}{主要資料}
\begingroup
% reader-facing Technical Notes break policy
\widowpenalty=10000
\clubpenalty=10000
\displaywidowpenalty=10000
\begin{tabularx}{\linewidth}{@{}>{\bfseries}p{0.22\linewidth}X@{}}
組織 & Microsoft \\
種別 & オープンウェイト \\
時系列 & 2024-04-23, 2024-05-21 \\
\end{tabularx}
\smallskip
\begin{minipage}{\linewidth}
% reader-facing Technical Notes event-only fact/source tail group
{\bfseries 一次資料から整理したtechnical points}
\begin{itemize}[leftmargin=1.5em,itemsep=0.35em]
\item \textbf{一次情報で確認できる事実}: Microsoftの一次資料により、Phi-3 familyについて2024-04-23にモデル公開、2024-05-21にモデル Family Expansionを確認した。 対象event近傍の一次資料から vision encoder を確認できる。
\end{itemize}
\begin{minipage}{\linewidth}
% reader-facing Technical Notes generic boundary/source tail group
\begin{samepage}
{\bfseries 一次資料}
\begingroup\sloppy
\Urlmuskip=0mu plus 2mu\relax
\begin{itemize}[leftmargin=1.5em,itemsep=0.25em]
\item {\scriptsize\url{https://azure.microsoft.com/en-us/blog/new-models-added-to-the-phi-3-family-available-on-microsoft-azure/}}
\item {\scriptsize\url{https://www.microsoft.com/en-us/research/publication/phi-3-technical-report-a-highly-capable-language-model-locally-on-your-phone/}}
\end{itemize}
\endgroup
\end{samepage}
\end{minipage}
\end{minipage}
\endgroup
\end{technicalnote}
\medskip

\Needspace{0.38\textheight}
\begin{technicalnote}{MobileLLM}{主要資料}
\begingroup
% reader-facing Technical Notes break policy
\widowpenalty=10000
\clubpenalty=10000
\displaywidowpenalty=10000
\begin{tabularx}{\linewidth}{@{}>{\bfseries}p{0.22\linewidth}X@{}}
組織 & - \\
種別 & 論文 \\
時系列 & 年表対象日付なし \\
\end{tabularx}
\smallskip
\begin{minipage}{\linewidth}
% reader-facing Technical Notes event-only fact/source tail group
{\bfseries 一次資料から整理したtechnical points}
\begin{itemize}[leftmargin=1.5em,itemsep=0.35em]
\item \textbf{一次情報で確認できる事実}: 選定済み一次資料では、MobileLLMはsub-billion parameter帯のon-device用途を対象に、単純な縮小ではなくdeep-and-thinなTransformer構成を中心に探索した。論文ではembedding sharingとgrouped-query attentionなど、parameter budgetを本体blockへ振り向ける設計を組み合わせ、125M/350M規模を含む小型モデルで既存同規模モデルとの比較を行う。精度・効率の数値は著者報告として扱い、実端末上の有用性はモデル品質だけでなくmemory/latency制約にも依存する。
\end{itemize}
\begin{minipage}{\linewidth}
% reader-facing Technical Notes generic boundary/source tail group
\begin{samepage}
{\bfseries 一次資料}
\begingroup\sloppy
\Urlmuskip=0mu plus 2mu\relax
\begin{itemize}[leftmargin=1.5em,itemsep=0.25em]
\item {\scriptsize\url{http://arxiv.org/abs/2402.14905v2}}
\end{itemize}
\endgroup
\end{samepage}
\end{minipage}
\end{minipage}
\endgroup
\end{technicalnote}
\medskip

\Needspace{0.38\textheight}
\begin{technicalnote}{OpenELM}{補足資料}
\begingroup
% reader-facing Technical Notes break policy
\widowpenalty=10000
\clubpenalty=10000
\displaywidowpenalty=10000
\begin{tabularx}{\linewidth}{@{}>{\bfseries}p{0.22\linewidth}X@{}}
組織 & Apple \\
種別 & オープンウェイト \\
時系列 & 2024-04-22 \\
\end{tabularx}
\smallskip
\begin{minipage}{\linewidth}
% reader-facing Technical Notes event-only fact/source tail group
{\bfseries 一次資料から整理したtechnical points}
\begin{itemize}[leftmargin=1.5em,itemsep=0.35em]
\item \textbf{一次情報で確認できる事実}: Appleの一次資料により、OpenELMについて2024-04-22にPaper First Submission、2024-04-22に研究公開を確認した。 対象event近傍の一次資料から image generation / small model / cost-efficient deployment を確認できる。
\end{itemize}
\begin{minipage}{\linewidth}
% reader-facing Technical Notes generic boundary/source tail group
\begin{samepage}
{\bfseries 一次資料}
\begingroup\sloppy
\Urlmuskip=0mu plus 2mu\relax
\begin{itemize}[leftmargin=1.5em,itemsep=0.25em]
\item {\scriptsize\url{http://arxiv.org/abs/2404.14619v2}}
\item {\scriptsize\url{https://machinelearning.apple.com/research/openelm}}
\end{itemize}
\endgroup
\end{samepage}
\end{minipage}
\end{minipage}
\endgroup
\end{technicalnote}
\medskip

\Needspace{0.38\textheight}
\begin{technicalnote}{Gemma family (Gemma / PaliGemma / Gemma 2)}{補足資料}
\begingroup
% reader-facing Technical Notes break policy
\widowpenalty=10000
\clubpenalty=10000
\displaywidowpenalty=10000
\begin{tabularx}{\linewidth}{@{}>{\bfseries}p{0.22\linewidth}X@{}}
組織 & Google \\
種別 & オープンウェイト \\
時系列 & 2024-02-21, 2024-05-14, 2024-06-27 \\
\end{tabularx}
\smallskip
\begin{minipage}{\linewidth}
% reader-facing Technical Notes event-only fact/source tail group
{\bfseries 一次資料から整理したtechnical points}
\begin{itemize}[leftmargin=1.5em,itemsep=0.35em]
\item \textbf{一次情報で確認できる事実}: Googleの一次資料により、Gemma family (Gemma / PaliGemma / Gemma 2)について2024-02-21にモデル公開、2024-05-14にモデル公開、2024-06-27にモデル公開を確認した。 対象event近傍の一次資料から 2B parameter scale / 7B parameter scale / 9B parameter scale / 27B parameter scale を確認できる。
\end{itemize}
\begin{minipage}{\linewidth}
% reader-facing Technical Notes generic boundary/source tail group
\begin{samepage}
{\bfseries 一次資料}
\begingroup\sloppy
\Urlmuskip=0mu plus 2mu\relax
\begin{itemize}[leftmargin=1.5em,itemsep=0.25em]
\item {\scriptsize\url{https://blog.google/innovation-and-ai/technology/developers-tools/gemma-open-models/}}
\item {\scriptsize\url{https://blog.google/innovation-and-ai/technology/developers-tools/google-gemma-2/}}
\item {\scriptsize\url{https://developers.googleblog.com/en/gemma-family-and-toolkit-expansion-io-2024/}}
\end{itemize}
\endgroup
\end{samepage}
\end{minipage}
\end{minipage}
\endgroup
\end{technicalnote}
\medskip

\Needspace{0.38\textheight}
\begin{technicalnote}{BitNet b1.58}{補足資料}
\begingroup
% reader-facing Technical Notes break policy
\widowpenalty=10000
\clubpenalty=10000
\displaywidowpenalty=10000
\begin{tabularx}{\linewidth}{@{}>{\bfseries}p{0.22\linewidth}X@{}}
組織 & - \\
種別 & 論文 \\
時系列 & 年表対象日付なし \\
\end{tabularx}
\smallskip
\begin{minipage}{\linewidth}
% reader-facing Technical Notes event-only fact/source tail group
{\bfseries 一次資料から整理したtechnical points}
\begin{itemize}[leftmargin=1.5em,itemsep=0.35em]
\item \textbf{一次情報で確認できる事実}: 選定済み一次資料では、BitNet b1.58はTransformerのLinear層のweightsを\{-1, 0, 1\}のternary値で表し、weight 1個あたり約1.58 bitという表現を学習時から前提にするBitNet変種である。論文は同parameter scaleのfull-precision Transformerとのperplexity・end-task比較と、乗算・memory footprint・energy面の効率化を報告しているが、数値性能とhardware効率は著者の評価条件に依存する。
\end{itemize}
\begin{minipage}{\linewidth}
% reader-facing Technical Notes generic boundary/source tail group
\begin{samepage}
{\bfseries 一次資料}
\begingroup\sloppy
\Urlmuskip=0mu plus 2mu\relax
\begin{itemize}[leftmargin=1.5em,itemsep=0.25em]
\item {\scriptsize\url{http://arxiv.org/abs/2402.17764v1}}
\end{itemize}
\endgroup
\end{samepage}
\end{minipage}
\end{minipage}
\endgroup
\end{technicalnote}
\medskip
